Fix admissible \(\rho\) and \(\delta\). Let \(X_i=x_0\) almost surely, and let \((G_i,S_i)\) be uniform on
\[
 \{2,3,\infty\}\times\{0,1\}.
\]
Conditional on \((G_i,S_i)=(a,s)\), generate \((A_i,B_i)\in\{-1,1\}^2\) according to
\[
 Q_{\rho,\delta}(A_i=u,B_i=v\mid G_i=a,S_i=s)
 =\frac{1+r_a uv}{4},
 \qquad u,v\in\{-1,1\},
\]
where \(r_2=r_3=0\) and \(r_\infty=\rho\). This is a probability distribution because \(|\rho|<\sqrt7/3<1\). It satisfies
\[
 E[A_i\mid G_i=a,S_i=s]=E[B_i\mid G_i=a,S_i=s]=0,
\]
\[
 E[A_i^2\mid G_i=a,S_i=s]=E[B_i^2\mid G_i=a,S_i=s]=1,
 \qquad
 E[A_iB_i\mid G_i=a,S_i=s]=r_a.
\]
Set
\[
 Z_{i1}=0,
 \qquad Z_{i2}=\frac{A_i}{8},
 \qquad Z_{i3}=\frac{B_i}{8}.
\]
For \(a\in\{2,3,\infty\}\), \(s\in\{0,1\}\), and \(t\in\{1,2,3\}\), define
\[
 Y_{it}(a,s)
 =Z_{it}+\delta s\mathbf 1\{a<\infty\}\mathbf 1\{t\ge a\},
 \qquad
 Y_{it}=Y_{it}(G_i,S_i).
\]
These specifications determine the full-data law \(Q_{\rho,\delta}\).

Consistency holds by construction. If \(t<a\), then \(Y_{it}(a,1)=Z_{it}=Y_{it}(\infty,1)\), proving no anticipation. For \(s=0\), \(Y_{it}(a,0)=Z_{it}=Y_{it}(\infty,0)\), proving no spillover to ineligibles. Moreover, under never enabling both eligibility states have outcome history \((Z_{i1},Z_{i2},Z_{i3})\). Hence every untreated eligibility contrast in every retained pairwise conditional DDD restriction is zero, so all such restrictions hold.

The retained cells are the six elements of \(\{2,3,\infty\}\times\{0,1\}\), so \(L_{\mathrm{ret}}=6\). Since \(X_i=x_0\),
\[
 p_{a,s}(x_0)=\frac16,
 \qquad
 0<\varepsilon=\frac18<\frac16=L_{\mathrm{ret}}^{-1},
 \qquad
 p_{a,s}(x_0)\ge\varepsilon.
\]
Thus overlap holds. Also
\[
 |Y_{it}(a,s)|\le\frac18+|\delta|<\frac14,
\]
and therefore \(E_P[\max_t|Y_{it}|^4]<1=C_4\).

Write \(k_1=(2,2)\), \(k_2=(2,3)\), and order the scores as
\[
 j_1=(k_1,3),
 \qquad j_2=(k_1,\infty),
 \qquad j_3=(k_2,\infty).
\]
The cell regressions are
\[
 m_{a,s,k_1}(x_0)=\delta s\mathbf 1\{a=2\},
 \qquad
 m_{a,s,k_2}(x_0)=\delta s\mathbf 1\{a\in\{2,3\}\}.
\]
Consequently
\[
 d_{k_1,3}(x_0)=d_{k_1,\infty}(x_0)
 =d_{k_2,\infty}(x_0)=\delta,
 \qquad
 \Theta_{k_1}(P)=\Theta_{k_2}(P)=\delta.
\]
Thus every target-cell contrast term in the centered scores vanishes. The residuals are
\[
 \Delta Y_{i,k_1}-m_{a,s,k_1}(x_0)=\frac{A_i}{8},
 \qquad
 \Delta Y_{i,k_2}-m_{a,s,k_2}(x_0)=\frac{B_i}{8}.
\]
Because \(p_{2,1}=\pi_{2,1}=p_{a,s}=1/6\), every nonzero residual loading is \(6\sigma_{a,s}^{2,c}\).

Each score contains four mutually exclusive cells, each with probability \(1/6\), squared loading \(36\), and residual variance \(1/64\). Therefore every diagonal covariance is \(3/8\). Scores \(j_1\) and \(j_2\) share exactly the two cohort-\(2\) cells with equal signs, giving covariance \(3/16\). Scores \(j_1\) and \(j_3\) share cohort \(2\), where \(E[A_iB_i\mid G_i=2,S_i]=0\), giving covariance zero. Scores \(j_2\) and \(j_3\) have zero covariance on cohort \(2\) and covariance \(\rho/64\) on each of the two never-enabled cells, giving covariance \(3\rho/16\). Lemma~\ref{lem:shared-covariance-expansion} thus yields
\[
 \Omega(\rho)=
 \begin{pmatrix}
 3/8&3/16&0\\
 3/16&3/8&3\rho/16\\
 0&3\rho/16&3/8
 \end{pmatrix}
 =\frac3{16}
 \begin{pmatrix}
 2&1&0\\
 1&2&\rho\\
 0&\rho&2
 \end{pmatrix}.
\]
The eigenvalues of the last unscaled matrix are
\[
 2,
 \qquad 2-\sqrt{1+\rho^2},
 \qquad 2+\sqrt{1+\rho^2},
\]
because its characteristic polynomial is \((2-\lambda)\{(2-\lambda)^2-(1+\rho^2)\}\). Since \(\rho^2<7/9\), \(\sqrt{1+\rho^2}<4/3\), and hence
\[
 \lambda_{\min}\{\Omega(\rho)\}>\frac18,
 \qquad
 \lambda_{\max}\{\Omega(\rho)\}<\frac58<8.
\]
Therefore \(8^{-1}I_3\preceq\Omega(\rho)\preceq8I_3\). Every defining condition of \(\mathcal M\) has now been verified, so \(Q_{\rho,\delta}\in\mathcal M\).

The first covariance block is \((3/16)\begin{pmatrix}2&1\\1&2\end{pmatrix}\), whose normalized inverse loading is \((1/2,1/2)^{\prime}\); the second block is a singleton. Thus
\[
 b^{\mathrm{sep}}=
 \begin{pmatrix}\omega_1/2\\ \omega_1/2\\ \omega_2\end{pmatrix}.
\]
The joint constraints are \(b_1+b_2=\omega_1\) and \(b_3=\omega_2\), so every feasible vector is
\[
 b(q)=
 \begin{pmatrix}
 \omega_1/2+q\\
 \omega_1/2-q\\
 \omega_2
 \end{pmatrix}.
\]
With \(h=(1,-1,0)^{\prime}\), direct differentiation gives
\[
 \frac{d}{dq}\{b(q)^{\prime}\Omega b(q)\}
 =\frac38(2q-\rho\omega_2),
 \qquad
 \frac{d^2}{dq^2}\{b(q)^{\prime}\Omega b(q)\}=\frac34>0.
\]
The unique minimizer occurs at \(q=\rho\omega_2/2\), proving
\[
 b^{\mathrm{joint}}=
 \begin{pmatrix}
 \omega_1/2+\rho\omega_2/2\\
 \omega_1/2-\rho\omega_2/2\\
 \omega_2
 \end{pmatrix}.
\]
The difference is \(b^{\mathrm{sep}}-b^{\mathrm{joint}}=-(\rho\omega_2/2)h\), and \(h^{\prime}\Omega h=3/8\). Orthogonality to the constrained minimizer therefore gives
\[
 V_{\mathrm{sep}}-V_{\mathrm{joint}}
 =\frac{\rho^2\omega_2^2}{4}\frac38
 =\frac{3\rho^2\omega_2^2}{32}>0.
\]
Finally, every post-adoption eligible effect equals \(\delta\ne0\). The parameter set \(\{(\rho,\delta):0<|\rho|<\sqrt7/3,\ 0<|\delta|<1/8\}\) is open and two-dimensional, and the laws vary with both parameters. This proves the nondegenerate open family.